\scanpage{38}
% L1–3
On en déduirait des \uuline{corollaires} : dans les deux cas $M$ est Coh-Mac, et si on désigne par $A'$ son anneau d'homothéties, $A'$ est Coh. Mac, et $M\simeq\Omega_{A'}$.

% L4–6
\uline{Problème} C Supposons \inserted{$\operatorname{Ext}^{n}(k,M)\ne0$ i.e. $H^n(M)\ne0$.} $\operatorname{Ext}^{n-1}(k,M)=\operatorname{Ext}^{n+1}(k,M)=0$, (la dernière \unsure{condition|condit.} \unsure{équivalant|équivalente} à : $H^n(M)$ injectif). \inserted{ou : $N\to\operatorname{Ext}^n(N,M)$ exact} Est-il vrai que $A$ est Coh. Mac. et $M\simeq\Omega_A^{m}$

% L7–10
[On notera qu'on a un hom naturel $M\to\Omega_A^{m}$, il faut montrer que c'est un isomorphisme]. \uline{Ici encore}, on est ramené (par le dévissage à la Serre) à prouver que l'annulateur de $\mathfrak m$ dans $M$ \uline{et} $A$ \deleted{est} sont nuls.

% archive sheet number: 37
